% \chapter{Appendix for Chapter 2}

% \section[More General GNS representations]{More General GNS representations}
% \label{AppendixA}



% {\color{magenta}
% Furthermore, any pure quasifree state defines a unique (up to unitary equivalence) irreducible GNS representation of the algebra with fairly explicit constructions. The GNS construction associated to a quasifree state relies on a map defined in the space of \textit{real} solutions  $K:C_\mathbb{R}\rightarrow \mathscr{H}$, where $KC_\mathbb{R} + \ii K C_\mathbb{R}$ is dense in $\mathscr{H}$ and 
% \begin{equation}
%     \langle Kf|Kg\rangle = \frac{1}{2}H(f,g) + \frac{\ii}{2}E(f,g) = W(f,g).
% \end{equation}
% The Hilbert space $\mathscr{H}$ can then be used to build the Hilbert space $\mathcal{F}(\mathscr{H})$ which is the Hilbert space in the GNS construction.

% To build $\mathscr{H}$ and the map $K$, we first notice that the Wightman and Hadamard distributions are positive in $C_\mathbb{R}$:
% \begin{equation}
%     W(f,f) = \tfrac{1}{2}H(f,f) = \omega(\hat{\phi}(f)\hat{\phi}(f)) = \omega(\hat{\phi}(f)^2) \geq 0
% \end{equation}
% Thus, the following inequality holds
% \begin{equation}
%     |E(f,g)|^2 \leq H(f,f)H(g,g),
% \end{equation}
% which can be proven by applying the Cauchy-Schwartz inequality to the positive bidistribution $W$. The Hilbert space $\mathscr{H}$ can be explicitly built in the case where the state $\omega$ is pure, which happens if and only if the inequality~\eqref{} is saturated, in the sense that
% \begin{equation}
%     H(f,f) = \sup_{g\in C_\mathbb{R}}\frac{|E(f,g)|^2}{H(g,g)}.
% \end{equation}
% In this case, the fact that $E$ is non-degenerate then implies that $H$ is non-degenerate. That is, $H$ defines a real positive definite inner product in $C_\mathbb{R}$, 
% \begin{equation}
%     \langle Kf|Kg\rangle_\mathbb{R} \coloneqq \frac{1}{2}H(f,g),\\    
% \end{equation}
% We can then define a real Hilbert space by taking $\mathscr{H}_\mathbb{R} = \overline{C_\mathbb{R}}$, the completion of $C_\mathbb{R}$ with respect to the norm induced by $\langle\cdot|\cdot\rangle_\mathbb{R}$. The operator $K$ is simply defined as the inclusion of $C_\mathbb{R}$ in its completion $\mathscr{H}_\mathbb{R}$. 

% Our next step is to extend $\mathscr{H}_\mathbb{R}$ to a complex Hilbert space. We can do this by noticing that the causal propagator, as a bilinear operator $E:C_\mathbb{R}\times C_\mathbb{R}\rightarrow \mathbb{R}$ is densely defined in $\mathcal{H}_\mathbb{R}$. The Riesz representation theorem ensures then that there exists a linear operator $J:\mathscr{H}_\mathbb{R}\rightarrow\mathscr{H}_\mathbb{R}$ such that
% \begin{equation}
%     \frac{1}{2}E(f,g) = \langle Kf| J Kg\rangle_\mathbb{R}.
% \end{equation}
% Antisymmetry of $E$ combined with the saturated inequality~\eqref{} imply that $J^\dagger = -J$ and $||J g||_\mathbb{R} = ||g||_\mathbb{R}$ for all $g\in\mathscr{H}_\mathbb{R}$. These properties then give that $J^2 = - \openone$, so that $J$ can be used to complexify the Hilbert space $\mathscr{H}_\mathbb{R}$, defining the multiplication by a complex scalar as $(a+\ii b)\cdot f \coloneqq a f - b J f$. Finally, we define the complexification of $\mathscr{H}_\mathbb{R}$ with the inner product
% \begin{equation}
%     \langle \psi_1|\psi_2\rangle \coloneqq \langle \psi_1| \psi_2\rangle_\mathbb{R} + \ii \langle \psi_1| J \psi_2\rangle_\mathbb{R}
% \end{equation}
% as the complex Hilbert space $\mathscr{H}$. This inner product turns out to satisfy the condition $\langle Kf|Kg\rangle = W(f,g)$. Indeed, if $\psi_1 = Kf$ and $\psi_2 = Kg$, then
% \begin{equation}
%     \langle Kf|Kg\rangle = \langle Kf| Kg\rangle_\mathbb{R} + \ii \langle Kf| J Kg\rangle_\mathbb{R} = \frac{1}{2}H(f,g) + \frac{\ii}{2}E(f,g) = W(f,g).
% \end{equation}



% It turns out that states in the space $\mathscr{H}$ can be fully parametrized by functions in $C_\mathbb{C}$ noticing that for real functions $f_\tc{r},f_\tc{i},g_\tc{r},g_\tc{i}\in C_\mathbb{R}$, $W(f_\tc{r}\! -\! \ii f_\tc{i},g_\tc{r} \!+ \!\ii g_\tc{i}) = \langle f_\tc{r}\!-\!J f_\tc{i}|g_\tc{r} \!-\! J g_\tc{i}\rangle$. That is, the space $\mathcal{H}$ is isomorphic to $C_\mathbb{C}$ with the inner product $\langle f|g\rangle = W(f^*\!,g)$. You may then wonder why we made such a complicated construction to achieve a result that could be state in one line. The reason for this is that each function $f\in C_\mathbb{C}$ canonically gives rise to a \textit{complex} solution of the Klein-Gordon equation $Ef$, which is certainly not what this construction entails, as it only uses the space of real solutions. The fact that the Hilbert space $\mathscr{H}$ happens to be isomorphic to $C_\mathbb{C}$ is almost accidental, due to the fact that we have complexified a Hilbert space built from real solutions.

% }


% \chapter{An Explicit Example Of Particle Detection}

% {\color{red}
% We choose the function $f(\mf x)$ by directly picking $\tilde{f}(-|\bm k|,-\bm k)$, which makes the process simpler due to the quotient in $C_{\mathbb{C}}$. For convenience we will pick $\tilde{f}(-|\bm k|,-\bm k)$ with a Gaussian-like profile, centered at $\bm k = \bm k_0$ with width $\delta$:
% \begin{equation}
%     \tilde{f}(-|\bm k|,- \bm k) = \frac{2\sqrt{2} (2\pi)^\frac{3}{2} |\bm k_0|^\frac{1}{2} |\bm k|e^{-\frac{|\bm k - \bm k_0|^2}{2\delta^2}}}{\delta^\frac{3}{2}\sqrt{2\pi |\bm k_0|\delta e^{-\frac{|\bm k_0|^2}{\delta^2}} + \pi^\frac{3}{2} (\delta^2 + 2 |\bm k_0|^2)\text{erf}(|\bm k_0|/\delta))}}.
% \end{equation}
% The normalization factor is picked to ensure normalization of the state $\ket{f}$, as discussed in~\eqref{}. The state $\ket{f}$ then corresponds to the positive mode classical solution 
% \begin{equation}
%     \phi(\mf x) = \int \dd^3 \bm k \, u_{\bm k}^*(f)u_{\bm k}(\mf x) = \frac{1}{(2\pi)^3}\int \frac{\dd^3 \bm k}{2|\bm k|}f(-|\bm k|,-\bm k) e^{\ii \mf k \cdot \mf x}
% \end{equation}
% }



\chapter{Evaluation of Smeared Field Bi-Distributions on Gaussian Spacetime Functions}\label{app:analytical}



In this appendix we will compute the bi-distributions $W(f_1,f_2)$, $E(f_1,f_2)$ and $H(f_1,f_2)$ evaluated at functions $f_\tc{i}(\mf x)$ defined by parameters $(T_\tc{i},t_\tc{i},\Omega_\tc{i}, \sigma_\tc{i}, \bm L_\tc{i})$:
\begin{equation}
    f_\tc{i}(\mf x) = \frac{e^{-\frac{|\bm x - \bm L_\tc{i}|^2}{2 \sigma_\tc{i}^2}}}{(2\pi \sigma_\tc{i}^2)^{3/2}}e^{- \frac{(t-t_\tc{i})^2}{2T_\tc{i}^2}}e^{\ii \Omega_\tc{i} t},
\end{equation}
for general parameters  $(T_1,t_1,\Omega_1, \sigma_1, \bm L_1)$ and $(T_2,t_2,\Omega_2, \sigma_2, \bm L_2)$. We will also compute $G_R(f_1,f_2)$, $G_A(f_1,f_2)$, $\Delta(f_1,f_2)$, and $G_F(f_1,f_2)$ in the case $\sigma_1 = \sigma_2 = \sigma$. Finally, we will discuss how to generalize our results for the computation of two-point functions that involve the momentum of a real massless quantum field and more general time derivatives of the field.

We start from Eq.~\eqref{eq:W0Fourier}, and notice that for the functions
\begin{equation}
    F_\tc{i}(\bm x) = \frac{e^{-\frac{|\bm x - \bm L_\tc{i}|^2}{2 \sigma_\tc{i}^2}}}{(2\pi \sigma_\tc{i}^2)^{3/2}}, \quad \quad \chi_\tc{i}(t) = e^{- \frac{(t-t_\tc{i})^2}{2T_\tc{i}^2}}e^{\ii \Omega_\tc{i} t},
\end{equation}
we have the Fourier transforms
\begin{equation}
    \tilde{F}(\bm k) = e^{\ii \bm k \cdot \bm L_\tc{i} - \frac{\sigma_\tc{i}^2 |\bm k|^2}{2}}, \quad \quad \tilde{\chi}(\omega) = \sqrt{2 \pi} T_\tc{i} e^{-\ii (\omega-\Omega_\tc{i}) t_\tc{i} - \frac{(\omega-\Omega_\tc{i})^2 T_\tc{i}^2}{2}}.
\end{equation}
Plugging the results above into Eq.~\eqref{eq:W0Fourier}, we find{\footnotesize
\begin{align}
    W(f_1,f_2) &= \frac{1}{(2\pi)^2} \int \frac{\dd^3\bm k}{2 |\bm k|} e^{\ii \bm k \cdot \bm L_1 - \frac{\sigma_1^2 |\bm k|^2}{2}} e^{-\ii \bm k \cdot \bm L_2 - \frac{\sigma_2^2 |\bm k|^2}{2}} T_2 e^{\ii (|\bm k|+\Omega_2) t_2 - \frac{(|\bm k|+\Omega_2)^2 T_2^2}{2}} T_1 e^{-\ii (|\bm k|-\Omega_1) t_1 - \frac{(|\bm k|-\Omega_1)^2 T_1^2}{2}}\nonumber\\
    &= \frac{T_1T_2e^{\ii (\Omega_1 t_1+\Omega_2 t_2)}}{(2\pi)^2} \int \frac{\dd^3\bm k}{2 |\bm k|} e^{\ii \bm k \cdot (\bm L_1 - \bm L_2)} e^{-\ii |\bm k|( t_1-t_2)}e^{ - \frac{(\sigma_1^2+\sigma_2^2) |\bm k|^2}{2}}   e^{- \frac{(|\bm k|+\Omega_2)^2 T_2^2}{2} - \frac{(|\bm k|-\Omega_1)^2 T_1^2}{2}} \nonumber\\
    &= \frac{T_1T_2e^{\ii (\Omega_1 t_1+\Omega_2 t_2)}}{(2\pi)^2} e^{- \frac{\Omega_2^2 T_2^2}{2} - \frac{\Omega_1^2 T_1^2}{2}}\int \frac{\dd^3\bm k}{2 |\bm k|} e^{\ii \bm k \cdot (\bm L_1 - \bm L_2)} e^{-\ii |\bm k|( t_1-t_2)}e^{ - \frac{(\sigma_1^2+\sigma_2^2+T_1^2 + T_2^2) |\bm k|^2}{2}}    e^{- |\bm k| (\Omega_2 T_2^2-\Omega_1 T_1^2)}\nonumber
    \\
    &= \frac{T_1T_2e^{\ii (\Omega_1 t_1+\Omega_2 t_2)}}{2\pi} e^{- \frac{\Omega_2^2 T_2^2}{2} - \frac{\Omega_1^2 T_1^2}{2}}\int \frac{\dd |\bm k|}{2 |\bm k|} |\bm k|^2 2 \,\text{sinc}(|\bm k|  |\bm L|) e^{-\ii |\bm k|t_0}e^{ - \frac{(\sigma_1^2+\sigma_2^2+T_1^2 + T_2^2) |\bm k|^2}{2}}    e^{- |\bm k| (\Omega_2 T_2^2-\Omega_1 T_1^2)}\nonumber\\
    &= \frac{T_1T_2e^{\ii (\Omega_1 t_1+\Omega_2 t_2)}}{2\pi|\bm L|} e^{- \frac{\Omega_2^2 T_2^2}{2} - \frac{\Omega_1^2 T_1^2}{2}}\int \dd |\bm k|\,\text{sin}(|\bm k|  |\bm L|) e^{-\ii |\bm k|t_0}e^{ - \frac{(\sigma_1^2+\sigma_2^2+T_1^2 + T_2^2) |\bm k|^2}{2}}    e^{- |\bm k| (\Omega_2 T_2^2-\Omega_1 T_1^2)},\label{eq:WfinalInt}
\end{align}}
where $\bm L = \bm L_1 - \bm L_2$ and $t_0 = t_1-t_2$. Using the result
\begin{equation}
    \int_0^\infty \dd r e^{- \ii b r}e^{- \frac{a^2r^2}{2}} = \frac{\sqrt{\pi}}{\sqrt{2}a}e^{- \frac{b^2}{2a^2}}\Big(1 - \ii \text{erfi}\left(\frac{b}{\sqrt{2}a}\right)\Big),
\end{equation}
and writing $\sin(|\bm k||\bm L|)$ as exponentials we find{\footnotesize
\begin{align}
     W(f_1,f_2) &= \frac{T_1T_2e^{\ii (\Omega_1 t_1+\Omega_2 t_2)}}{2\pi|\bm L|}  \frac{1}{2\ii}\frac{\sqrt{\pi}e^{- \frac{\Omega_2^2 T_2^2}{2} - \frac{\Omega_1^2 T_1^2}{2}}}{\sqrt{2} \sqrt{T_1^2 + T_2^2 + \sigma_1^2 + \sigma_2^2}}\Bigg(e^{- \frac{(t_0 - |\bm L|+\ii (\Omega_1 T_1^2- \Omega_2 T_2^2))^2}{2(T_1^2 + T_2^2 + \sigma_1^2 + \sigma_2^2)}}\Bigg(1 - \ii \text{erfi}\left(\tfrac{t_0 - |\bm L|+\ii (\Omega_1 T_1^2- \Omega_2 T_2^2)}{\sqrt{2}\sqrt{T_1^2 + T_2^2 + \sigma_1^2 + \sigma_2^2}}\right)\Bigg)\nonumber\\
     &\:\:\:\:\:\:\:\:\:\:\:\:\:\:\:\:\:\:\:\:\:\:\:\:\:\:\:\:\:\:\:\:\:\:\:\:\:\:\:\:\:\:\:\:\:\:\:\:\:\:\:\:\:\:\:\:\:\:\:\:\:\:\:\:\:\:\:\:\:\:\:\:\:\:\:\:\:\:\:\:\:\:\:\:\:\:\:\:\:\:\:\:\:-e^{- \frac{(t_0 + |\bm L|+\ii (\Omega_1 T_1^2- \Omega_2 T_2^2))^2}{2(T_1^2 + T_2^2 + \sigma_1^2 + \sigma_2^2)}}\Bigg(1 - \ii \text{erfi}\left(\tfrac{t_0 + |\bm L|+\ii (\Omega_1 T_1^2- \Omega_2 T_2^2)}{\sqrt{2}\sqrt{T_1^2 + T_2^2 + \sigma_1^2 + \sigma_2^2}}\right)\Bigg)\nonumber\\
     &= \frac{T_1T_2e^{\ii (\Omega_1 t_1+\Omega_2 t_2)}e^{- \frac{\Omega_2^2 T_2^2}{2} - \frac{\Omega_1^2 T_1^2}{2}}}{4\sqrt{2\pi}|\bm L|\sqrt{T_1^2 + T_2^2 + \sigma_1^2 + \sigma_2^2}} \Bigg(e^{- \frac{(t_0 - |\bm L|+\ii (\Omega_1 T_1^2- \Omega_2 T_2^2))^2}{2(T_1^2 + T_2^2 + \sigma_1^2 + \sigma_2^2)}}\Bigg(-\ii - \text{erfi}\left(\tfrac{t_0 - |\bm L|+\ii (\Omega_1 T_1^2- \Omega_2 T_2^2)}{\sqrt{2}\sqrt{T_1^2 + T_2^2 + \sigma_1^2 + \sigma_2^2}}\right)\Bigg)\nonumber\\
     &\:\:\:\:\:\:\:\:\:\:\:\:\:\:\:\:\:\:\:\:\:\:\:\:\:\:\:\:\:\:\:\:\:\:\:\:\:\:\:\:\:\:\:\:\:\:\:\:\:\:\:\:\:\:\:\:\:\:\:\:\:\:\:\:\:\:\:\:\:\:\:\:\:\:\:\:\:\:\:+e^{- \frac{(t_0 + |\bm L|+\ii (\Omega_1 T_1^2- \Omega_2 T_2^2))^2}{2(T_1^2 + T_2^2 + \sigma_1^2 + \sigma_2^2)}}\Bigg(\ii + \text{erfi}\left(\tfrac{t_0 + |\bm L|+\ii (\Omega_1 T_1^2- \Omega_2 T_2^2)}{\sqrt{2}\sqrt{T_1^2 + T_2^2 + \sigma_1^2 + \sigma_2^2}}\right)\Bigg)\Bigg).
\end{align}}
And, as it turns out, we find that{\footnotesize
\begin{align}
    H(f_1,f_2) &= \frac{T_1T_2e^{\ii (\Omega_1 t_1+\Omega_2 t_2)}e^{- \frac{\Omega_2^2 T_2^2}{2} - \frac{\Omega_1^2 T_1^2}{2}}}{2\sqrt{2\pi}|\bm L|\sqrt{T_1^2 + T_2^2 + \sigma_1^2 + \sigma_2^2}} \Bigg(e^{- \frac{(t_0 - |\bm L|+\ii (\Omega_1 T_1^2- \Omega_2 T_2^2))^2}{2(T_1^2 + T_2^2 + \sigma_1^2 + \sigma_2^2)}}\text{erfi}\left(\frac{|\bm L|-t_0-\ii (\Omega_1 T_1^2- \Omega_2 T_2^2)}{\sqrt{2}\sqrt{T_1^2 + T_2^2 + \sigma_1^2 + \sigma_2^2}}\right)\\
     &\:\:\:\:\:\:\:\:\:\:\:\:\:\:\:\:\:\:\:\:\:\:\:\:\:\:\:\:\:\:\:\:\:\:\:\:\:\:\:\:\:\:\:\:\:\:\:\:\:\:\:\:\:\:\:\:\:\:\:\:\:\:\:\:\:\:\:\:\:\:\:\:\:\:\:\:\:\:\:+e^{- \frac{(t_0 + |\bm L|+\ii (\Omega_1 T_1^2- \Omega_2 T_2^2))^2}{2(T_1^2 + T_2^2 + \sigma_1^2 + \sigma_2^2)}}\text{erfi}\left(\frac{|\bm L|+t_0+\ii (\Omega_1 T_1^2- \Omega_2 T_2^2)}{\sqrt{2}\sqrt{T_1^2 + T_2^2 + \sigma_1^2 + \sigma_2^2}}\right)\Bigg),\label{eq:Hbig}\nonumber
\end{align}}
and
\begin{align}
    E(f_1,f_2) &= \frac{T_1T_2e^{\ii (\Omega_1 t_1+\Omega_2 t_2)}e^{- \frac{\Omega_2^2 T_2^2}{2} - \frac{\Omega_1^2 T_1^2}{2}}}{2\sqrt{2\pi}|\bm L|\sqrt{T_1^2 + T_2^2 + \sigma_1^2 + \sigma_2^2}} \Bigg(e^{- \frac{(t_0 + |\bm L|+\ii (\Omega_1 T_1^2- \Omega_2 T_2^2))^2}{2(T_1^2 + T_2^2 + \sigma_1^2 + \sigma_2^2)}}-e^{- \frac{(t_0 - |\bm L|+\ii (\Omega_1 T_1^2- \Omega_2 T_2^2))^2}{2(T_1^2 + T_2^2 + \sigma_1^2 + \sigma_2^2)}}\Bigg).
\end{align}

We now move on to the Green's functions, which will allow us to compute the Feynman propagator and the symmetric propagator. In order to compute these, we will resort to integration in spacetime of the spacetime smearing functions, using the expressions for the Green's functions of Eqs.~\eqref{eq:GRGAdelta}. Unfortunately, we will have to consider a more restricted parameter space where $\sigma_1 = \sigma_2 = \sigma$ in order to solve these integrals. For the retarded Green's function, we have{\footnotesize
\begin{align}
    G_R(f_1,f_2) &= \int \dd V\dd V' f_1(\mf x) f_2(\mf x') G_R(\mf x,\mf x')\\
    & = \frac{1}{(2\pi)^3 \sigma^6} \int \dd t \dd t' \dd^3\bm x \dd^3 \bm x' e^{- \frac{|\bm x - \bm L_1|^2}{2\sigma^2}}e^{- \frac{|\bm x' - \bm L_2|^2}{2\sigma^2}} e^{\ii \Omega_1 t} e^{- \frac{(t-t_1)^2}{2T_1^2}}e^{\ii \Omega_2 t'} e^{- \frac{(t'-t_2)^2}{2T_2^2}}\left(- \frac{\delta(t' - t + |\bm x - \bm x'|)}{4\pi|\bm x-\bm x'|} \right)\nonumber\\
    & = -\frac{1}{2(2\pi)^4 \sigma^6} \int \dd t  \dd^3\bm x \dd^3 \bm x' e^{- \frac{|\bm x - \bm L_1|^2}{2\sigma^2}}e^{- \frac{|\bm x' - \bm L_2|^2}{2\sigma^2}} e^{\ii \Omega_1 t} e^{- \frac{(t-t_1)^2}{2T_1^2}}e^{\ii \Omega_2 (t - |\bm x - \bm x'|)} e^{- \frac{(t - |\bm x - \bm x'|-t_2)^2}{2T_2^2}}\frac{1}{|\bm x - \bm x'|}.\nonumber
\end{align}}
We now perform the change of variables
\begin{align}
    \bm x = \frac{1}{\sqrt{2}}(\bm u +\bm v) + \bm L_1, \quad\quad
    \bm x' = \frac{1}{\sqrt{2}}(\bm u -\bm v) + \bm L_1,
\end{align}
so that $|\bm x - \bm x'| = \sqrt{2} |\bm v|$. Defining $\bm L = \bm L_1 - \bm L_2$, the integral becomes:{\footnotesize
\begin{align}
    G_R(f_1,f_2) & = -\frac{1}{2(2\pi)^4 \sigma^6} \int \dd^3\bm u \dd^3 \bm v e^{- \frac{|\bm u + \bm v|^2}{4\sigma^2}}e^{- \frac{|\bm u - \bm v + \sqrt{2}\bm L|^2}{4\sigma^2}}\frac{1}{\sqrt{2}|\bm v|} \int \dd t e^{\ii \Omega_1 t} e^{- \frac{(t-t_1)^2}{2T_1^2}}e^{\ii \Omega_2 (t - \sqrt{2}|\bm v|)} e^{- \frac{(t - \sqrt{2}|\bm v|-t_2)^2}{2T_2^2}}\\
    & = -\frac{1}{2(2\pi)^4 \sigma^6} \int   \dd^3\bm u \dd^3 \bm v e^{- \frac{|\bm u|^2 + |\bm v|^2}{4\sigma^2}}e^{- \frac{|\bm u|^2 +|\bm v|^2 + 2|\bm L|^2}{4\sigma^2}}e^{-\frac{\sqrt{2}\bm u\cdot \bm L + \sqrt{2}\bm v \cdot \bm L}{2\sigma^2}} \frac{1}{\sqrt{2}|\bm v|} \nonumber\\
    &\:\:\:\:\:\:\:\:\:\:\:\:\:\:\:\:\:\:\:\:\:\:\:\:\:\:\:\:\:\:\:\:\:\:\:\:\:\:\:\:\:\:\:\:\:\:\:\:\:\:\:\:\:\:\:\:\:\:\:\:\:\:\:\:\:\:\:\:\:\:\:\:\:\:\:\:\:\:\:\:\:\:\:\:\:\:\:\:\:\:\:\:\:\:\:\:\:\:\:\:\:\:\:\:\:\times\int \dd t e^{\ii \Omega_1 t} e^{- \frac{(t-t_1)^2}{2T_1^2}}e^{\ii \Omega_2 (t - \sqrt{2}|\bm v|)} e^{- \frac{(t - \sqrt{2}|\bm v|-t_2)^2}{2T_2^2}}\nonumber\\
    &= -\frac{e^{- \frac{|\bm L|^2}{2\sigma^2}}}{2(2\pi)^4 \sigma^6} \int   \dd^3\bm u \dd^3 \bm v e^{-\frac{|\bm u|^2}{2\sigma^2}}e^{-\frac{|\bm v|^2}{2\sigma^2}} e^{-\frac{\bm u\cdot \bm L}{\sqrt{2}\sigma^2}}e^{\frac{\bm v \cdot \bm L}{\sqrt{2}\sigma^2}}\frac{1}{\sqrt{2}|\bm v|} \int \dd t e^{\ii \Omega_1 t} e^{- \frac{(t-t_1)^2}{2T_1^2}}e^{\ii \Omega_2 (t - \sqrt{2}|\bm v|)} e^{- \frac{(t - \sqrt{2}|\bm v|-t_2)^2}{2T_2^2}}.\nonumber
\end{align}}
We can now proceed to first solve for the angular integrals in $\bm u$ and $\bm v$, and then the radial integral in $|\bm u|$, resulting in{\footnotesize
\begin{align}
    G_R(f_1,f_2) & = -\frac{e^{- \frac{|\bm L|^2}{2\sigma^2}}}{2(2\pi)^2 \sigma^6} \int   \dd |\bm u| \dd |\bm v| \,|\bm u|^2 |\bm v|^2 e^{-\frac{|\bm u|^2}{2\sigma^2}}e^{-\frac{|\bm v|^2}{2\sigma^2}} \frac{2\sqrt{2} \sigma^2\sinh(\frac{|\bm L||\bm u|}{\sqrt{2}\sigma^2})}{|\bm u||\bm L|}\frac{2\sqrt{2} \sigma^2\sinh(\frac{|\bm L||\bm v|}{\sqrt{2}\sigma^2})}{|\bm v||\bm L|}\frac{1}{\sqrt{2}|\bm v|}\nonumber \\
    &\:\:\:\:\:\:\:\:\:\:\:\:\:\:\:\:\:\:\:\:\:\:\:\:\:\:\:\:\:\:\:\:\:\:\:\:\:\:\:\:\:\:\:\:\:\:\:\:\:\:\:\:\:\:\:\:\:\:\:\:\:\:\:\:\:\:\:\:\:\:\:\:\:\:\:\:\:\:\:\:\:\:\:\:\:\:\:\:\:\:\:\:\:\:\:\:\:\:\:\:\:\:\:\:\:\times\int \dd t e^{\ii \Omega_1 t} e^{- \frac{(t-t_1)^2}{2T_1^2}}e^{\ii \Omega_2 (t - \sqrt{2}|\bm v|)} e^{- \frac{(t - \sqrt{2}|\bm v|-t_2)^2}{2T_2^2}}\nonumber\\
    & = -\frac{e^{- \frac{|\bm L|^2}{2\sigma^2}}}{\sqrt{2}\pi^2\sigma^2|\bm L|^2} \int  \dd |\bm u| \dd |\bm v| \,|\bm u| e^{-\frac{|\bm u|^2}{2\sigma^2}} \sinh(\frac{|\bm L||\bm u|}{\sqrt{2}\sigma^2})e^{-\frac{|\bm v|^2}{2\sigma^2}}\sinh(\frac{|\bm L||\bm v|}{\sqrt{2}\sigma^2})\nonumber\\
    &\:\:\:\:\:\:\:\:\:\:\:\:\:\:\:\:\:\:\:\:\:\:\:\:\:\:\:\:\:\:\:\:\:\:\:\:\:\:\:\:\:\:\:\:\:\:\:\:\:\:\:\:\:\:\:\:\:\:\:\:\:\:\:\:\:\:\:\:\:\:\:\:\:\:\:\:\:\:\:\:\:\:\:\:\:\:\:\:\:\:\:\:\:\:\:\:\:\:\:\:\:\:\:\:\:\times \int \dd t e^{\ii \Omega_1 t} e^{- \frac{(t-t_1)^2}{2T_1^2}}e^{\ii \Omega_2 (t - \sqrt{2}|\bm v|)} e^{- \frac{(t - \sqrt{2}|\bm v|-t_2)^2}{2T_2^2}}\nonumber\\
    & = -\frac{e^{- \frac{|\bm L|^2}{2\sigma^2}}}{\sqrt{2}\pi^2\sigma^2|\bm L|^2} \int  \dd |\bm v| \, \frac{\sqrt{\pi}\sigma |\bm L|}{2}e^{\frac{|\bm L|^2}{4 \sigma^2}}e^{-\frac{|\bm v|^2}{2\sigma^2}} \sinh(\frac{|\bm L||\bm v|}{\sqrt{2}\sigma^2}) \int \dd te^{\ii \Omega_1 t} e^{- \frac{(t-t_1)^2}{2T_1^2}}e^{\ii \Omega_2 (t - \sqrt{2}|\bm v|)} e^{- \frac{(t - \sqrt{2}|\bm v|-t_2)^2}{2T_2^2}}\nonumber\\
    & = -\frac{e^{- \frac{|\bm L|^2}{4\sigma^2}}}{2\sqrt{2}\pi^{3/2}\sigma|\bm L|} \int \dd |\bm v| \, e^{-\frac{|\bm v|^2}{2\sigma^2}} \sinh(\frac{|\bm L||\bm v|}{\sqrt{2}\sigma^2}) \int \dd t e^{\ii \Omega_1 t} e^{- \frac{(t-t_1)^2}{2T_1^2}}e^{\ii \Omega_2 (t - \sqrt{2}|\bm v|)} e^{- \frac{(t - \sqrt{2}|\bm v|-t_2)^2}{2T_2^2}}.\label{eq:GRfinalInt}
\end{align}}
Defining $t_0 = t_1 - t_2$, the integral in $t$ yields{\footnotesize
\begin{align}
     \int \dd t e^{\ii \Omega_1 t} e^{- \frac{(t-t_1)^2}{2T_1^2}}e^{\ii \Omega_2 (t - \sqrt{2}|\bm v|)}& e^{- \frac{(t - \sqrt{2}|\bm v|-t_2)^2}{2T_2^2}}\nonumber \textcolor{white}{\text{If you noticed the mistake, congratulations! The final result for GA is correct :)}}\\
     &\!\!\!\! = \frac{\sqrt{2\pi}T_1 T_2}{\sqrt{T_1^2 + T_2^2}}e^{- \frac{|\bm v|^2}{T_1^2 + T_2^2}}e^{\frac{\sqrt{2}|\bm v|(t_0- \ii (\Omega_1 T_1^2 - \Omega_2 T_2^2))}{(T_1^2 + T_2^2)}}e^{-\frac{t_0^2}{2(T_1^2 + T_2^2)}-\frac{T_1^2T_2^2(\Omega_1+\Omega_2)^2}{2(T_1^2 + T_2^2)}+\ii(\Omega_1 + \Omega_2)\frac{(t_1 T_2^2 + t_2 T_1^2)}{T_1^2 + T_2^2}}.
\end{align}}
The final integral in $|\bm v|$ is merely a combination of Gaussian integrals. After simplifications, we finally find{\footnotesize
\begin{align}
    G_R(f_1,f_2) = - \frac{T_1T_2e^{\ii (\Omega_1 t_1+\Omega_2 t_2)}e^{- \frac{\Omega_2^2 T_2^2}{2} - \frac{\Omega_1^2 T_1^2}{2}}}{4\sqrt{2\pi}|\bm L|\sqrt{T_1^2 + T_2^2 + 2\sigma^2}}\Bigg(&e^{- \frac{(t_0 - |\bm L|+\ii (\Omega_1 T_1^2- \Omega_2 T_2^2))^2}{2(T_1^2 + T_2^2 + 2\sigma^2)}}\left(1 + \text{erf}\left(\tfrac{|\bm L|(T_1^2+T_2^2)+2\sigma^2(t_0+\ii (\Omega_1 T_1^2- \Omega_2 T_2^2))}{2\sigma\sqrt{T_1^2 + T_2^2}\sqrt{T_1^2 + T_2^2 + 2\sigma^2}}\right)\right)\nonumber\\
    &\!\!\!\!\!\!\!\!\!\!\!\!\!\!\!\!\!\!+e^{- \frac{(t_0 + |\bm L|+\ii (\Omega_1 T_1^2- \Omega_2 T_2^2))^2}{2(T_1^2 + T_2^2 + 2\sigma^2)}}\left(-1+\text{erf}\left(\tfrac{|\bm L|(T_1^2+T_2^2)-2\sigma^2(t_0+\ii (\Omega_1 T_1^2- \Omega_2 T_2^2))}{2\sigma\sqrt{T_1^2 + T_2^2}\sqrt{T_1^2 + T_2^2 + 2\sigma^2}}\right)\right)\Bigg).\label{eq:GRbig}
\end{align}}
Combining $G_R(f_1,f_2)$ and $E(f_1,f_2)$, we can then find $G_A(f_1,f_2)$,{\footnotesize
\begin{align}
    G_A(f_1,f_2) = - \frac{T_1T_2e^{\ii (\Omega_1 t_1+\Omega_2 t_2)}e^{- \frac{\Omega_2^2 T_2^2}{2} - \frac{\Omega_1^2 T_1^2}{2}}}{4\sqrt{2\pi}|\bm L|\sqrt{T_1^2 + T_2^2 + 2\sigma^2}}\Bigg(&e^{- \frac{(t_0 - |\bm L|+\ii (\Omega_1 T_1^2- \Omega_2 T_2^2))^2}{2(T_1^2 + T_2^2 + 2\sigma^2)}}\left(-1 + \text{erf}\left(\tfrac{|\bm L|(T_1^2+T_2^2)+2\sigma^2(t_0+\ii (\Omega_1 T_1^2- \Omega_2 T_2^2))}{2\sigma\sqrt{T_1^2 + T_2^2}\sqrt{T_1^2 + T_2^2 + 2\sigma^2}}\right)\right)\nonumber\\
    &\!\!\!\!\!\!\!\!\!\!\!\!\!\!\!\!\!\!+e^{- \frac{(t_0 + |\bm L|+\ii (\Omega_1 T_1^2- \Omega_2 T_2^2))^2}{2(T_1^2 + T_2^2 + 2\sigma^2)}}\left(1+\text{erf}\left(\tfrac{|\bm L|(T_1^2+T_2^2)-2\sigma^2(t_0+\ii (\Omega_1 T_1^2- \Omega_2 T_2^2))}{2\sigma\sqrt{T_1^2 + T_2^2}\sqrt{T_1^2 + T_2^2 + 2\sigma^2}}\right)\right)\Bigg).\label{eq:GAbig}
\end{align}}
Finally, we find $\Delta(f_1,f_2)$ by adding Eqs.~\eqref{eq:GRbig} and~\eqref{eq:GAbig}:{\footnotesize
\begin{align}
    \Delta(f_1,f_2) = - \frac{T_1T_2e^{\ii (\Omega_1 t_1+\Omega_2 t_2)}e^{- \frac{\Omega_2^2 T_2^2}{2} - \frac{\Omega_1^2 T_1^2}{2}}}{2\sqrt{2\pi}|\bm L|\sqrt{T_1^2 + T_2^2 + 2\sigma^2}}\Bigg(&e^{- \frac{(t_0 - |\bm L|+\ii (\Omega_1 T_1^2- \Omega_2 T_2^2))^2}{2(T_1^2 + T_2^2 + 2\sigma^2)}} \text{erf}\left(\tfrac{|\bm L|(T_1^2+T_2^2)+2\sigma^2(t_0+\ii (\Omega_1 T_1^2- \Omega_2 T_2^2))}{2\sigma\sqrt{T_1^2 + T_2^2}\sqrt{T_1^2 + T_2^2 + 2\sigma^2}}\right)\nonumber\\
    &+e^{- \frac{(t_0 + |\bm L|+\ii (\Omega_1 T_1^2- \Omega_2 T_2^2))^2}{2(T_1^2 + T_2^2 + 2\sigma^2)}}\text{erf}\left(\tfrac{|\bm L|(T_1^2+T_2^2) - 2\sigma^2(t_0+\ii (\Omega_1 T_1^2- \Omega_2 T_2^2))}{2\sigma\sqrt{T_1^2 + T_2^2}\sqrt{T_1^2 + T_2^2 + 2\sigma^2}}\right)\Bigg).\label{eq:Deltabig}
\end{align}}
Using Eqs.~\eqref{eq:Hbig} and~\eqref{eq:Deltabig}, one can then find $G_F(f_1,f_2) = \frac{1}{2}H(f_1,f_2) + \frac{\ii}{2}\Delta(f_1,f_2)$.


In order to obtain the field bi-distributions evaluated at the field's momentum $\hat{\pi}(\mf x) = \partial_t \hat{\phi}(\mf x)$ rather than at the field's amplitude smeared in Gaussian spacetime regions, one can simply differentiate the results of this section with respect to the parameters $\Omega_1$, $\Omega_2$, using
\begin{equation}
    \omega( \hat{\pi}(f_1)\hat{\phi}(f_2)) = \omega( \hat{\phi}(-\partial_tf_1)\hat{\phi}(f_2)),
\end{equation}
and noticing that for the functions used here
\begin{equation}
    \partial_t f_1(\mf x) = \left(- \frac{t}{T_1^2} + \frac{t_1}{T_1^2} + \ii \Omega_1\right) f_1(\mf x) = \frac{\ii}{T_1^2}\dv{}{\Omega_1}f_1(\mf x) + \left(\frac{t_1}{T_1^2} + \ii \Omega_1\right) f_1(\mf x).
\end{equation}
That is, one finds that
\begin{equation}
    \omega( \hat{\pi}(f_1)\hat{\phi}(f_2)) =  - \frac{\ii}{T_1^2}\dv{}{\Omega_1}W(f_1,f_2) -\left(\frac{t_1}{T_1^2} + \ii \Omega_1\right) W(f_1,f_2).
\end{equation}
Analogous expressions are valid for higher derivatives of the field in both arguments, and for the other bi-distributions $H$, $E$, $G_R$, $G_A$, $G_F$ and $\Delta$. For brevity, we will not write these explicit expressions here, but they can be straightforwardly computed.



\chapter{Estimation of the Fermi bound}\label{app:fermi}

In this appendix we present an estimation for the Fermi bound, that is, we find an approximate bound for the maximum radius that a region contained in the rest spaces $\Sigma_\tau$ can have in order to be contained within the normal neighbourhood of the point $\mf z(\tau)$.

In order to obtain our estimate, we use the expansion of Eq. \eqref{eq:expansionFNC} for the
metric in Fermi normal coordinates. It is important to notice that the approximations of Eq. \eqref{eq:expansionFNC} are not only to second order in the distance of points to the curve, but are also the first order in curvature and acceleration expansions of the metric in Fermi normal coordinates. In the regime where these expansions are valid, the Fermi bound can be estimated as the maximum radius such that the metric does not become degenerate. That is, the largest radius such that the metric of Eq. \eqref{eq:expansionFNC}
is invertible. A good estimate for the $\tau$-Fermi bound can be given by considering the largest value of $|\bm x|$ such that 
\begin{equation}
    -(1+a_i x^i) - R_{0i0j}x^ix^j \neq 0.
\end{equation}
Notice that at $\bm x = 0$ this gives $-1$, so that for each $\tau$ we are looking for a bound on the radius $\ell_r$ such that 
\begin{equation}\label{eq:condition}
    |\bm x| \leq \ell_r \Rightarrow -(1+a_i x^i) - R_{0i0j}x^ix^j < 0. 
\end{equation}

In order to find $\ell_r$, let $a = \sqrt{a_ia^i}$ and define
\begin{equation}
    \lambda_R = \max_{|\bm x| = 1}\left(-R_{0i0j}x^ix^j\right).
\end{equation}
That is, $\lambda_R$ corresponds to the largest negative eigenvalue of $R_{0i0j}$, if there are any, else it is zero. We will now show that
\begin{equation}
    \ell_r = \frac{1}{a+\sqrt{\lambda_R}}
\end{equation}
fulfills the condition of Eq. \eqref{eq:condition}. Let $|\bm x|<\ell_r$, then
\begin{align}
    |\bm x|(a + \sqrt{\lambda_R})<1 \:\:\:\Rightarrow\:\:\: 1 - a|\bm x| >\sqrt{\lambda_R}|\bm x|,
\end{align}
so that using $a_i x^i \geq -a |\bm x|$,
\begin{equation}\label{eq:noname}
    1 + a_i x^i> \sqrt{\lambda_R} |\bm x| \:\:\:\Rightarrow\:\:\: (1 + a_i x^i)^2 > \lambda_R |\bm x| \geq - R_{0i0j}x^i x^j.
\end{equation}
This gives the desired result $0>-(1 + a_i x^i)^2- R_{0i0j}x^i x^j$.

We then obtain the approximate lower bound for the $\tau$-Fermi bound,
\begin{equation}\label{eq:tauFermiApprox}
    \ell_\tau \gtrsim \frac{1}{a+\sqrt{\lambda_R}},
\end{equation}
which is valid for points $\mf x$ such that $|\bm x|$ is sufficiently smaller than the curvature radius of spacetime and $1/a$, which is the regime of validity of the approximation of Eq. \eqref{eq:expansionFNC}. Notice that this bound works exactly in the case where there is no spacetime curvature, giving \mbox{$\ell_r = 1/a$}. The bound for the Fermi bound is then obtained by taking the infimum of Eq. \eqref{eq:tauFermiApprox} with respect to $\tau$.


%\chapter{Equivalence of the formalism for different choices of Fermi coordinates around $\mf{z}(\tau)$}

%In this section we show that the framework presented in Subsection \ref{sub:main} is independent of the orthonormal frame $\mf{e}_i(\tau_0)$ upon which the Fermi normal coordinates are built. For this purpose, it is enough to show that for any other orthonormal frame $\mf{e}_{i'}(\tau_0)$, one would obtain the same operators $\hat{\bm x}$ and $\hat{\bm p}$, acting on $L^2(\Sigma_{\tau})$ for every $\tau$. At $\tau_0$, the two orthonormal frames $\mf{e}_i(\tau_0)$ and $\mf e_{i'}(\tau_0)$ can be related by an orthogonal operator $R^{i'}{}_i$, such that $\mf{e}_i(\tau_0) = R^{i'}{}_i\mf e_{i'}(\tau_0)$.


\chapter{The Redshift Factor and the Metric Determinant}\label{app:det}

The metric in Fermi normal coordinates can be written as
\begin{equation}\label{eq:ADM}
    g = g_{\tau\tau}\dd\tau^2 +2g_{\tau i} \dd \tau \dd x^i+h_{ij}\dd x^i \dd x^j.
\end{equation}
Then, the inverse metric reads
\begin{equation}
    g^{-1} = g^{\tau\tau}\partial_\tau\otimes\partial_\tau + g^{\tau i}\partial_\tau \otimes \partial_i+g^{i\tau}\partial_i \otimes \partial_\tau + g^{ij}\partial_i\otimes\partial_j,
\end{equation}
where
\begin{align}
    g^{\tau\tau} &= \frac{1}{g_{\tau\tau} - g_{\tau i}g_{\tau j} h^{ij}},\\
    g^{\tau i} &= -\frac{h^{ij}g_{\tau j}}{g_{\tau\tau} - g_{\tau k}g_{\tau l} h^{kl}},\\
    g^{ij} &= h^{ij} + \frac{h^{ik}g_{\tau k}h^{jl}g_{\tau l}}{g_{\tau\tau} - g_{\tau k}g_{\tau l} h^{kl}},
\end{align}
and $h^{ij}$ is the inverse of the spatial metric, satisfying $h^{ik}h_{kj} = \delta^i_j$. The metric determinant can be written as
\begin{equation}
    \sqrt{-g} = \sqrt{h^{ij}g_{\tau i} g_{\tau j} -g_{\tau\tau}} \sqrt{g_{\Sigma}},
\end{equation}
where $g_{\Sigma} = \det(h_{ij})$.


The redshift factor associated to the foliation \mbox{$\tau = \text{const.}$} is determined by the norm of the 1-form $\dd \tau$, which is given by 
\begin{equation}
    {|\dd \tau|}^2 = g^{\tau\tau} = \frac{1}{g_{\tau\tau} - g_{\tau i}g_{\tau j} h^{ij}}.
\end{equation}
The redshift factor is then given by
\begin{equation}
    \gamma(\mf x) = \frac{1}{|\dd\tau|} = \abs{g_{\tau\tau} - g_{\tau i}g_{\tau j} h^{ij}}^\frac{1}{2},
\end{equation}
so that the invariant volume element of spacetime can be written as
\begin{equation}
    \dd V = \frac{\dd \tau}{|\dd \tau|}\wedge \dd \Sigma = \dd\tau \wedge(\gamma(\mf x) \dd \Sigma),
\end{equation}
where $\dd \Sigma = \sqrt{g_{\Sigma}}\,\dd^n \bm x$ is the measure in the \mbox{$\tau = \text{const}.$} surfaces.



\definecolor{DeepBlue}{RGB}{30,50,100}


\chapter{The Localization Condition for Quantum Fields}\label{app:locCond}

We can show that  Eq.~\eqref{eq:locCond} holds for a localized field in Minkowski spacetime defined by the Lagrangian~\eqref{eq:Lagphid}. The Riemann normal coordinates in Minkowski spacetime are simply inertial coordinates, and we can write the class of functions ${f}_{\mf x_0}$ simply as $f_{\mf x_0}(t,\bm x) = f(t-t_0,\bm x-\bm x_0)$ in the inertial coordinates $(t,\bm x)$ where the potential $V(\bm x)$ is defined. We will essentially show Eq.~\eqref{eq:locCond} when this restriction is imposed over the shape regions. In this case, the limit~\eqref{eq:locCond} corresponds to the limit $|\bm x_0|\to \infty$. 

Notice that for any compactly smooth real function $h$, we have
\begin{equation}\label{eq:auxlocCond}
    \bra{0_\tc{d}}\hat{\phi}_\tc{d}(f_{\mf x_0})\hat{\phi}_\tc{d}(h)\ket{0_\tc{d}} = \sum_{\bm n} u_{\bm n}({f}_{\mf x_0})u_{\bm n}^*(h),
\end{equation}
where the sum is absolutely convergent, as it corresponds to the inner product between the one-particle states $\ket{f_{\mf x_0}}$ and $\ket{h}$ (see~\eqref{eq:Wfg1part}).  We then have
\begin{equation}
    u_{\bm n}(f_{\mf x_0}) = \int \dd t \,\dd^3 \bm x \, f(t-t_0,\bm x - \bm x_0) \Phi_{\bm n}(\bm x) e^{-\ii \omega_{\bm n} t} = \int \dd t \,\dd^3 \bm x \, f(t,\bm x) \Phi_{\bm n}(\bm x + \bm x_0) e^{-\ii \omega_{\bm n} (t+t_0)}. 
\end{equation}
Due to the fact that the potential $V(\bm x)$ goes to infinity, we also have that the asymptotic behaviour of the functions $\Phi_{\bm n}(\bm x)$ at large $\bm x$ is at most exponential~\cite{expDecay}, so that we can write
 \begin{equation}
     \int \dd^3 \bm x f(t,\bm x) \Phi_{\bm n}(\bm x + \bm x_0) = \chi_{\bm n}(t,\bm x_0) e^{-\gamma_{\bm n}|\bm x_0|}, 
 \end{equation}
 with $\gamma_{\bm n}> \gamma^*>0$ for a given $\gamma^*$ and
 \begin{equation}
     \lim_{|\bm x_0| \to \infty}\chi_{\bm n}(t,\bm x_0) e^{-\gamma^*|\bm x_0|} = 0.
 \end{equation}
 We then find
\begin{equation}
    u_{\bm n}(f_{\mf x_0}) = \int \dd t \chi_{\bm n}(t, \bm x_0)e^{-\ii \omega_{\bm n} (t+t_0)}e^{-\gamma_{\bm n}|\bm x_0|} = \tilde{\chi}_{\bm n}(\omega_{\bm n},\bm x_0)e^{-\ii \omega_{\bm n} t_0}e^{-\gamma_{\bm n}|\bm x_0|},
\end{equation}
where $\tilde{\chi}_{\bm n}$ denotes the Fourier transform. Also notice that the eigenvalues of $L = - \nabla^2 + m_\tc{d}^2 + V(\bm x)$ with positive $V(\bm x)$ constitute an unbounded positive sequence, so that $\omega_{\bm n}\to \infty$ as $\bm n \to \infty$, also implying that $\tilde{\chi}(\omega_{\bm n},\bm x_0)\to 0$ for each $\bm x_0$, using the Riemann-Lebesgue Lemma.

We can now take the limit $|\bm x_0|\to \infty$ in Eq.~\eqref{eq:auxlocCond} using the dominated convergence theorem, as we have a product of the bounded terms $u_{\bm n}^*(h)$, $\tilde{\chi}_{\bm n}(\omega_{\bm n},\bm x_0)e^{-\ii \omega_{\bm n} t_0}e^{-\gamma_{\bm n}|\bm x_0|}\leq \tilde{\chi}_{\bm n}(\omega_{\bm n},\bm x_0)e^{-\ii \omega_{\bm n} t_0}e^{-\gamma^*|\bm x_0|}$ in an absolutely convergent series:
\begin{equation}
    \lim_{|\bm x_0|\to \infty }\bra{0_\tc{d}}\hat{\phi}_\tc{d}(f_{\mf x_0})\hat{\phi}_\tc{d}(h)\ket{0_\tc{d}} = \sum_{\bm n} \lim_{|\bm x_0|\to \infty} u_{\bm n}^*(h)\tilde{\chi}_{\bm n}(\omega_{\bm n},\bm x_0) e^{-\ii \omega_{\bm n} t_0} e^{-\gamma^*|\bm x_0|} = 0.
\end{equation}
The result above also holds when $h = f_{\mf x_0}$ with an analogous proof, so that we have shown that
\begin{equation}
    \lim_{|\bm x_0|\to \infty }\bra{0_\tc{d}}\hat{\phi}_\tc{d}(f_{\mf x_0})^2\ket{0_\tc{d}} =  0.
\end{equation}
Now notice that due to the vacuum $\ket{0_\tc{d}}$ being quasifree, all other terms of the form $\bra{0_\tc{d}}\hat{\phi}_\tc{d}(f_{\mf x_0})^n\ket{0_\tc{d}}$ either vanish if $n$ is odd or are proportional to $\bra{0_\tc{d}}\hat{\phi}_\tc{d}(f_{\mf x_0})^2\ket{0_\tc{d}}$, thus showing that the state $\ket{0_\tc{d}}$ is localized with this restriction of shape functions.

Moreover, given any quasifree state in $\mathcal{F}(\mathscr{H}_\tc{d})$, the expected values of operators of the form $\hat{\phi}_\tc{d}(f_{\mf x_0})^2$ will be given by sums and products of $W(h,{f}_{\mf x_0})$, implying that any quasifree state in the GNS representation of $\ket{0_\tc{d}}$ is also a localized state, from Eq.~\eqref{eq:Wfg1part}. For a Gaussian state $\omega_\tc{g}$ that is not quasifree, we can use the fact that $\omega_\tc{g}(\hat{\phi}_\tc{d}(f_{\mf x_0}))^2 \leq \omega_\tc{g}(\hat{\phi}_\tc{d}(f_{\mf x_0})^2)$, which also implies that the limit of $|\bm x_0|\to \infty$ of $\omega_\tc{g}(\hat{\phi}(f_{\mf x_0})^n)$ vanishes for all $n$.

Finally, notice that although we do not provide an explicit generalization of this result in general globally hyperbolic spacetimes, it is clear that the result still holds, provided that the mode functions $\Phi_{\bm n}(\bm x)$ are exponentially bounded by the proper distance of $\mf x_0$ and any other event $\mf y$ in the same Cauchy surface, which is achieved by a wide class of localizing potentials.


\chapter{Corrections to the Zeeman Interaction to Leading Order in the Fine Structure Constant}\label{app:derivation}


The goal of this appendix is to obtain the result of Eq.~\eqref{hflatzeemanexact}, which shows corrections to the Zeeman effect that depend on the shape of the radial function of the electron field modes corresponding to the electron state. This amounts to computing the integral of the function $\upphi(r)$, defined in Eq.~\eqref{eq:phir}, in space.

In Eq.~\eqref{psi-sol-atom} we showed the form of the $s$ orbital modes of an electron in terms of the radial functions $f(r)$ and $g(r)$. When $j=1/2$ and $p=+1$ these radial functions satisfy the differential equations~\cite{RQMgreiner}
\begin{align}
    \dv{g}{r} - (E-V(r)+m_e)f(r) &= 0,\\
    \dv{f}{r} +\frac{2}{r} f(r) + (E-V(r)-m_e)g(r) &= 0,
\end{align}
where $E$ is the energy level of the given orbital. The differential equations above are valid for any central potential $V(r)$. We can obtain an expression for $\upphi(r)$ by multiplying the first equation by $g(r)$ and the second equation by $f(r)$, yielding
\begin{align}
    \frac{1}{2}\dv{g^2}{r} - (E-V(r)+m_e)f(r)g(r) &= 0,\\
    \frac{1}{2}\dv{f^2}{r} +\frac{2}{r} f(r)g(r) + (E-V(r)-m_e)f(r)g(r) &= 0.
\end{align}
Adding the two equations we find
\begin{align}
    2m_e f(r) g(r) = \frac{1}{2}\dv{}{r}\left(f^2(r) + g^2(r)\right)+ \frac{2}{r} f^2(r).
\end{align}
Integrating the result above from $\infty$ to $r$ we find
\begin{equation}
    \upphi(r) = \int_{\infty}^r \dd r' f(r') g(r') = \frac{1}{4m_e}(f^2(r) + g^2(r)) + \frac{1}{m_e} \int_\infty^r \dd r' \frac{f^2(r')}{r'},
\end{equation}
where we directly integrated the total derivative using that $f(\infty) = g(\infty) = 0$. 

The integral of $\upphi(r)$ in space is then given by
\begin{equation}
    \int \dd^3 \bm x \, \upphi(|\bm x|) = \frac{4\pi}{4m_e} \int_0^\infty \dd r \, r^2 (f^2(r) + g^2(r)) - \frac{4\pi}{m_e}\int_0^\infty \dd r  \int_r^\infty  \dd r'\frac{r^2}{r'} f^2(r'),
\end{equation}
where we picked up factors of $4\pi$ due to integration over the angular coordinates and spherical symmetry of the functions involved. Normalization of the spinor spherical harmonics and of the mode functions $\psi_{\bm N}(\bm x)$ implies that
\begin{equation}\label{eq:normfg}
    \int_0^\infty \dd r \, r^2 (f^2(r) + g^2(r)) = 1.
\end{equation}
To handle the double integral, we can reparametrize the region defined by $0<r<\infty$ and $r<r'<\infty$ as $0<r'<\infty$ and $0<r<r'$. This gives
\begin{equation}\label{eq:changerrp}
    \int_0^\infty \dd r  \int_r^\infty  \dd r'\frac{r^2}{r'} f^2(r') = \int_0^\infty \dd r'  \int_0^{r'}  \dd r\frac{r^2}{r'} f^2(r') = \frac{1}{3}\int_0^\infty \dd r' (r')^2 f^2(r') = \frac{1}{3}\int_0^\infty \dd r\, r^2 f^2(r).
\end{equation}
Combining the results of Eqs.~\eqref{eq:normfg} and~\eqref{eq:changerrp}, we find
\begin{equation}
    \int \dd^3 \bm x \, \upphi(|\bm x|) = \frac{\pi}{m_e} - \frac{4\pi}{3m_e}\int_0^\infty \dd r\, r^2 f^2(r).
\end{equation}
Thus, with a constant magnetic field, we find
\begin{equation}
    \hat{H}_I(t) = - \frac{q}{2\pi} \int d^3\bm x\, \upphi(r)\, \hat{\bm{\sigma}}\cdot\bm{B}(t) = - \frac{q}{2m_e}\left(1 - \frac{4}{3}\int_0^\infty \dd r\, r^2 f^2(r)\right)\hat{\bm{\sigma}}\cdot\bm{B}(t).
\end{equation}

Importantly, this result holds true for modes with $j=1/2$ and $p=+1$ defined by any central potential $V(r)$. The corrections are then entirely determined by the function $f(r)$ and are of the order of the inverse of the product of the mass of the electron and the effective localization of the bound states.


{


\chapter{Two Localized Fields Interacting with a Klein-Gordon Field}\label{app:twoQFTs}

Denote the localized fields by $\hat{\phi}_\tc{a}$ and $\hat{\phi}_\tc{b}$, and the free field that they both interact with by $\hat{\phi}$ in the regions defined by the supports of $\zeta_\tc{a}(\mf x)$ and $\zeta_\tc{b}(\mf x)$, respectively. The interaction Hamiltonian density of the theory will be prescribed as
\begin{equation}
    \hat{h}_I(\mf x) = \lambda  (\zeta_\tc{a}(\mf x)\hat{\phi}(\mf x)\hat{\phi}_\tc{a}(\mf x) + \zeta_\tc{b}(\mf x)\phi(\mf x) \hat{\phi}_\tc{b}(\mf x)).
\end{equation}
We now write the $\phi_{\tc{a},\tc{b}}$ fields with the spacetime smearing functions as
\begin{align}
\zeta_\tc{a}(\mf x)\hat{\phi}_\tc{a}(\mf x) &= \sum_{\bm n} \zeta_\tc{a}(\mf x)\left( u_{\bm n}^{\tc{a}}(\mf x) \hat{a}_{\bm n}^\tc{a} +  u^{\tc{a}*}_{\bm n}(\mf x)\hat{a}^{\tc{a}\dagger}_{\bm n} \right) = \sum_{\bm n} \hat{Q}^{\tc{a}}_{\bm n}(\mf x)\\
\zeta_\tc{b}(\mf x)\hat{\phi}_\tc{b}(\mf x) &= \sum_{\bm n} \zeta_\tc{b}(\mf x)\left( u_{\bm n}^{\tc{b}}(\mf x) \hat{a}_{\bm n}^\tc{b} +  u^{\tc{b}*}_{\bm n}(\mf x) \hat{a}_{\bm n}^{\tc{b}\dagger} \right) = \sum_{\bm n}   \hat{Q}_{\bm n}^\tc{b}(\mf x),
\end{align}
where $u_{\bm n}^{\tc{a}}(\mf x) = e^{-\ii \omega_{\bm n}^{\tc{a}} t}\Phi^\tc{a}_{\bm n}(\bm x)$, $u_{\bm n}^{\tc{b}}(\mf x) = e^{-\ii \omega_{\bm n}^{\tc{b}} t}\Phi_{\bm n}^{\tc{b}}(\bm x)$, and the field expansion will depend on the specific boundary conditions and equations of motion. We are working under the assumption that the field has discrete energy levels, implying that the sums above are discrete. The field expansions automatically define states $\ket{0_\tc{a}}$ and $\ket{0_\tc{b}}$, which are annihilated by all operators $\hat{a}_{\bm n}^\tc{a}$ and $\hat{a}_{\bm n}^\tc{b}$, respectively. We can then write the Hamiltonian interaction density as
\begin{equation}
    \hat{h}_I(\mf x) = \lambda\hat{\phi}(\mf x) \sum_{\bm n} \left(\hat{Q}_{\bm n}^{\tc{a}}(\mf x) +  \hat{Q}_{\bm n}^\tc{b}(\mf x)\right).
\end{equation}
In perturbation theory, we then get
\begin{equation}
    \hat{U} = \mathcal{T}\exp(- \ii \int \dd V \hat{h}_I(\mf x)) = \openone + \hat{U}^{(1)} + \hat{U}^{(2)} + \mathcal{O}(\lambda^3),
\end{equation}
where:
\begin{align}
    \hat{U}^{(1)} = - \ii \int \dd V \hat{h}_I(\mf x) = - \ii \lambda \int \dd V  \hat{\phi}(\mf x)  \sum_{\bm n} \left(\hat{Q}^\tc{a}_{\bm n}(\mf x) +  Q_{\bm n}^\tc{b}(\mf x)\right).
\end{align}{\footnotesize
\begin{align}
    \hat{U}^{(2)}& =  - \int \dd V\dd V' \hat{h}_I(\mf x) \hat{h}_I(\mf x')\theta(t-t')\\&= - \int \dd V \dd V' \hat{\phi}(\mf x)\hat{\phi}(\mf x')\theta(t-t') \sum_{nm} \Big( \hat{Q}_{\bm n}^{\tc{a}}(\mf x)\hat{Q}_{\bm m}^{\tc{a}}(\mf x')+ \hat{Q}^\tc{b}_{\bm n}(\mf x)\hat{Q}_{\bm m}^{\tc{b}}(\mf x')+  \hat{Q}^\tc{a}_{\bm n}(\mf x)\hat{Q}^\tc{b}_{\bm m}(\mf x') +  \hat{Q}^\tc{b}_{\bm n}(\mf x) \hat{Q}^\tc{a}_{\bm m}(\mf x')\Big).
\end{align}}
The final state will be given by
\begin{equation}
    \hat{\rho}_f = \hat{U}\hat{\rho}_0\hat{U}^\dagger = \hat{\rho}_0 + \hat{U}^{(1)}\hat{\rho}_0 + \hat{\rho}_0 \hat{U}^{(1)\dagger} + \hat{U}^{(1)}\hat{\rho}_0\hat{U}^{(1)\dagger}+ \hat{U}^{(2)}\hat{\rho}_0 + \hat{\rho}_0 \hat{U}^{(2)\dagger} + \mathcal{O}(\lambda^3).
\end{equation}
We will assume that $\hat{\rho}_0 = \ket{0_\tc{a}0_\tc{b}}\!\!\bra{0_\tc{a}0_\tc{b}}\otimes \hat{\rho}_\phi = \hat{\rho}_{0,\tc{ab}}\otimes \hat{\rho}_\phi$ and that $\tr_\phi(\hat{U}^{(1)}\hat{\rho}_0) = 0$, so that the $\mathcal{O}(\lambda)$ terms do not contribute to the partial state of the cavities A and B. We then only need to compute \mbox{$\tr_\phi(\hat{U}^{(1)}\hat{\rho}_0\hat{U}^{(1)\dagger}+ \hat{U}^{(2)}\hat{\rho}_0 + \hat{\rho}_0 \hat{U}^{(2)\dagger})$}. We have:
{\footnotesize
\begin{align}
    &\tr_\phi(\hat{U}^{(1)}\hat{\rho}_0\hat{U}^{(1)\dagger}) \nonumber\\&= \lambda^2\int \dd V\dd V' W(\mf x', \mf x)  \sum_{nm} \Big( \hat{Q}^{\tc{a}}_{\bm n}(\mf x)\hat{\rho}_{0,\tc{ab}}\hat{Q}^{\tc{a}}_{\bm m}(\mf x')+ \hat{Q}^\tc{b}_{\bm n}(\mf x)\hat{\rho}_{0,\tc{ab}}\hat{Q}_{\bm m}^{\tc{b}}(\mf x')+ \hat{Q}^\tc{a}_{\bm n}(\mf x)\hat{\rho}_{0,\tc{ab}}\hat{Q}^{\tc{b}}_{\bm m}(\mf x') +  \hat{Q}^\tc{b}_{\bm n}(\mf x) \hat{\rho}_{0,\tc{ab}}\hat{Q}^{\tc{a}}_{\bm m}(\mf x')\Big),
\end{align}}
and{\footnotesize
\begin{align}
    \tr_\phi(\hat{U}^{(2)}\hat{\rho}_0) &= - \lambda^2\int \dd V \dd V' W(\mf x, \mf x')\theta(t-t')  \nonumber\\\times&\sum_{nm} \Big( \hat{Q}^{\tc{a}}_{\bm n}(\mf x)\hat{Q}^{\tc{a}}_{\bm m}(\mf x')\hat{\rho}_{0,\tc{ab}}+ \hat{Q}^\tc{b}_{\bm n}(\mf x)\hat{Q}_{\bm m}^{\tc{b}}(\mf x')\hat{\rho}_{0,\tc{ab}}+  \hat{Q}^\tc{a}_{\bm n}(\mf x)\hat{Q}^\tc{b}_{\bm m}(\mf x')\hat{\rho}_{0,\tc{ab}} +  \hat{Q}^\tc{b}_{\bm n}(\mf x) \hat{Q}^\tc{a}_{\bm m}(\mf x')\hat{\rho}_{0,\tc{ab}}\Big),
\end{align}}
where $W(\mf x, \mf x') = \text{tr}(\hat{\rho}_\phi \hat{\phi}(\mf x) \hat{\phi}(\mf x'))$. Notice that 
\begin{equation}
    \hat{\rho}_{0,\tc{ab}} = \ket{0_\tc{a}0_\tc{b}}\!\!\bra{0_\tc{a}0_\tc{b}} = \bigotimes_{nm} \ket{0_{\bm n}^\tc{a}0_{\bm m}^\tc{b}}\!\!\bra{0_{\bm n}^\tc{a}0_{\bm m}^\tc{b}} = \hat{\rho}_{\tc{d},0}\bigotimes_{n,m>1} \ket{0_{\bm n}^\tc{a}0_{\bm m}^\tc{b}}\!\!\bra{0_{\bm n}^\tc{a}0_{\bm m}^\tc{b}},
\end{equation}
where $\ket{0_{\bm n}^{\tc{a}}}$ and $\ket{0_{\bm n}^{\tc{b}}}$ denotes the ground states of each harmonic of each cavity, and 
\begin{equation}
    \hat{\rho}_{\tc{d},0} = \ket{0_1^\tc{a}0_1^\tc{b}}\!\!\bra{0_1^\tc{a}0_1^\tc{b}}
\end{equation}
is the ground state of the first harmonic in each cavity. We then trace out all cavity modes except for the first harmonic. That is, we will compute the density operator \mbox{$\tr_{\phi,H}(\hat{\rho}_f) = \tr_H(\tr_\phi(\hat{U}^{(1)}\hat{\rho}_0\hat{U}^{(1)\dagger}+ \hat{U}^{(2)}\hat{\rho}_0 + \hat{\rho}_0 \hat{U}^{(2)\dagger}))$}, where $\tr_H$ denotes the trace over all cavity harmonics, except the first, for fields A and B.

Overall, we will need to compute the trace of quantities of the form $\tr_H(\hat{Q}^\tc{a}_{\bm n}(\mf x)\hat{Q}^\tc{a}_{\bm m}(\mf x')\hat{\rho}_{0,\tc{ab}})$. We know that since $\hat{Q}^\tc{i}_{\bm n}(\mf x) = (u_{\bm n}^{\tc{i}}(\mf x)\hat{a}_{\bm n}^\tc{i} + u_{\bm n}^{\tc{i}*}(\mf x)\hat{a}^{\tc{i}\dagger}_{\bm n})$ for $\tc{I} = \tc{A},\tc{B}$. Therefore, the products of the form $\hat{Q}^\tc{i}_{\bm n}(\mf x)\hat{Q}_{\bm m}^\tc{i}(\mf x')$ will only give non-diagonal elements if $n = m$. When $n= m \neq 1$, we have:
\begin{equation}
    \tr_\tc{a}(\hat{Q}^\tc{a}_{\bm n}(\mf x)\hat{Q}^\tc{a}_{\bm n}(\mf x')\ket{0^\tc{a}_{\bm n}}\!\!\bra{0^\tc{a}_{\bm n}}) = \bra{0_{\bm n}^\tc{a}}\hat{Q}^\tc{a}_{\bm n}(\mf x)\hat{Q}^\tc{a}_{\bm n}(\mf x')\ket{0_{\bm n}^\tc{a}} = \zeta_\tc{a}(\mf x)\zeta_\tc{a}(\mf x')u_{\bm n}^\tc{a}(\mf x) u_{\bm n}^{\tc{a}*}(\mf x')
\end{equation}
and
\begin{equation}
    \tr_\tc{b}(\hat{Q}^\tc{b}_{\bm n}(\mf x)\hat{Q}^\tc{b}_{\bm n}(\mf x')\ket{0^\tc{b}_{\bm n}}\!\!\bra{0^\tc{b}_{\bm n}}) = \bra{0_{\bm n}^\tc{b}}\hat{Q}^\tc{b}_{\bm n}(\mf x)\hat{Q}^\tc{b}_{\bm n}(\mf x')\ket{0_{\bm n}^\tc{b}} = \zeta_\tc{b}(\mf x) \zeta_\tc{b}(\mf x')u_{\bm n}^\tc{b}(\mf x) u_{\bm n}^{\tc{b}*}(\mf x').
\end{equation}
We find that for $n$ and $m$ different from $1$,
\begin{equation}
    \tr_H(\hat{Q}_{\bm n}^\tc{i}(\mf x)\hat{Q}_{\bm m}^\tc{j}(\mf x')\hat{\rho}_{0,\tc{ab}}) = \delta_{\bm{nm}} \delta_{\tc{i}\tc{j}} \zeta_\tc{i}(\mf x)\zeta_\tc{j}(\mf x')u_{\bm n}^{\tc{i}}(\mf x)u_{\bm m}^{\tc{j}*}(\mf x')\hat{\rho}_{\tc{d},0},
\end{equation}
where $\tr_H(\hat{Q}_{\bm n}^\tc{a}(\mf x)\hat{Q}_{\bm m}^\tc{b}(\mf x')\hat{\rho}_{0,\tc{ab}}) = 0$ automatically, as it factors into expectation values of creation and annihilation operators in $\tc{A}$ and $\tc{B}$ evaluated at the vacuum. Finally, notice that when $n = m = 1$ we do not need to trace over it, because $H$ encompasses every harmonic except for the first one.

Putting the results above together, we then find that{\footnotesize
\begin{align}
    \tr_H(\tr_\phi(\hat{U}^{(1)} \hat{\rho}_0\hat{U}^{(1)\dagger})) = \!\lambda^2\!\!\!\int &\!\dd V \dd V'  W(\mf x' ,\mf x)\\
    & \times \Big( \hat{Q}^{\tc{a}}_1(\mf x)\hat{\rho}_{\tc{d},0}\hat{Q}_1^{\tc{a}}(\mf x')+ \hat{Q}_1^\tc{b}(\mf x)\hat{\rho}_{\tc{d},0}\hat{Q}_1^{\tc{b}}(\mf x')+ \hat{Q}^\tc{a}_1(\mf x)\hat{\rho}_{\tc{d},0}\hat{Q}^{\tc{b}}_1(\mf x') +  \hat{Q}_1^\tc{b}(\mf x) \hat{\rho}_{\tc{d},0}\hat{Q}^{\tc{a}}_1(\mf x') \nonumber\\&+\sum_{n>1}( \zeta_\tc{a}(\mf x)\zeta_\tc{a}(\mf x')u_{\bm n}^{\tc{a}}(\mf x) u_{\bm n}^{\tc{a}*}(\mf x')+\zeta_\tc{b}(\mf x)\zeta_\tc{b}(\mf x')u_{\bm n}^{\tc{b}}(\mf x) u_{\bm n}^{\tc{b}*}(\mf x'))\hat{\rho}_{\tc{d},0}\Big)\nonumber
\end{align}}
and{\footnotesize
\begin{align}
    \!\tr_H(\tr_\phi(\hat{U}^{(2)}\hat{\rho}_0))\! =\! - \lambda^2\!\!\!\int\! \!\dd V \dd V' &W(\mf x, \mf x')\theta(t-t')\\
    &\times\Big( \hat{Q}_1^{\tc{a}}(\mf x)\hat{Q}_1^{\tc{a}}(\mf x')\hat{\rho}_{\tc{d},0}\!+\! \hat{Q}^\tc{b}_1(\mf x)\hat{Q}_1^{\tc{b}}(\mf x')\hat{\rho}_{\tc{d},0}\!+\! \hat{Q}^\tc{a}_1(\mf x)\hat{Q}^\tc{b}_1(\mf x')\hat{\rho}_{\tc{d},0}\! +\! \hat{Q}^\tc{b}_1(\mf x) \hat{Q}^\tc{a}_1(\mf x')\hat{\rho}_{\tc{d},0}\nonumber\\
   &\:\:\:\:\:\:\:\:\:\:\:\:\:\:\:\:\:\:+\sum_{n>1} (\zeta_\tc{a}(\mf x)\zeta_\tc{a}(\mf x')u_{\bm n}^{\tc{a}}(\mf x)  u_{\bm n}^{\tc{a}*}(\mf x')+ \zeta_\tc{b}(\mf x)\zeta_\tc{b}(\mf x')u_{\bm n}^{\tc{b}}(\mf x) u_{\bm n}^{\tc{b}*}(\mf x'))\hat{\rho}_{\tc{d},0}\Big).\nonumber
   \end{align}}
The last term $\tr_H(\tr_\phi(\hat{\rho}_0 \hat{U}^{(2)\dagger}))$ is simply the conjugate of the term above. Notice that the terms proportional to $\hat{\rho}_{\tc{d},0}$ will cancel when all terms are added together. This can be seen from an explicit calculation using $\theta(t-t') + \theta(t'-t) = 1$, or simply by noticing that each term in the Dyson expansion is traceless due to trace preservation. %The final result is then exactly what is found when two harmonic oscillator UDW detectors to the field smeared by functions $u^\tc{a}_1(\bm x)$ and $u^\tc{b}_1(\bm x)$, respectively.

The products of terms $\hat{Q}_1^\tc{i}(\mf x)\hat{Q}_1^\tc{j}(\mf x')$ is given by
{\footnotesize
\begin{align}
    \hat{Q}_1^\tc{i}(\mf x)\hat{Q}_1^\tc{j}(\mf x') &= \zeta_{\tc{i}}(\mf x)\zeta_{\tc{j}}(\mf x') (u_1^\tc{i}(\mf x) \hat{a}_{\bm n}^\tc{i} + u_1^{\tc{i}*}(\mf x) \hat{a}_{\bm n}^{\tc{i}\dagger}) (u_1^\tc{j}(\mf x') \hat{a}_{\bm n}^\tc{j} + u_1^{\tc{j}*}(\mf x') \hat{a}_{\bm n}^{\tc{j}\dagger}) \\
    &= \zeta_{\tc{i}}(\mf x)\zeta_{\tc{j}}(\mf x') (u_1^\tc{i}(\mf x) u_1^\tc{j}(\mf x') \hat{a}_{\bm n}^\tc{i}\hat{a}_{\bm n}^\tc{j} + u_1^{\tc{i}*}(\mf x)u_1^\tc{j}(\mf x') \hat{a}_{\bm n}^{\tc{i}\dagger}\hat{a}_{\bm n}^\tc{j} + u_1^\tc{i}(\mf x)u_1^{\tc{j}*}(\mf x') \hat{a}_{\bm n}^\tc{i} \hat{a}_{\bm n}^{\tc{j}\dagger} + u_1^{\tc{i}*}(\mf x) u_1^{\tc{j}*}(\mf x') \hat{a}_{\bm n}^{\tc{i}\dagger}\hat{a}_{\bm n}^{\tc{j}\dagger}).\nonumber
\end{align}}
We then define the spacetime smearing functions
\begin{align}
    \Lambda_\tc{a}(\mf x) &\coloneqq \zeta_\tc{a}(\mf x) u^{\tc{a}}_1(\mf x),\\
    \Lambda_\tc{b}(\mf x) &\coloneqq \zeta_\tc{b}(\mf x) u^{\tc{b}}_1(\mf x),    
\end{align}
so that the $\hat{Q}_1^{\tc{i}}(\mf x)$ terms read simply as
\begin{align}
    \hat{Q}_1^{\tc{a}}(\mf x) = \Lambda_\tc{a}(\mf x) \hat{a}^{\tc{a}}_1+\Lambda^*_\tc{a}(\mf x) \hat{a}^{\tc{a}\dagger}_1,\\
    \hat{Q}_1^{\tc{b}}(\mf x) = \Lambda_\tc{b}(\mf x) \hat{a}^{\tc{b}}_1+\Lambda^*_\tc{b}(\mf x) \hat{a}^{\tc{b}\dagger}_1.
\end{align}
We then see that the final state of the fields $\tc{A}$ and $\tc{B}$ can be written as
\begin{align}
    \hat\rho_\tc{d} = \tr_{\phi,H}(\hat\rho_f) = \hat{\rho}_{\tc{d},0} + \tr_\phi\left(\hat{U}_I^{(1)} \hat{\rho}_{\tc{d},0}\hat{U}_I^{(1)^\dagger} + \hat{U}_I^{(2)} \hat{\rho}_{\tc{d},0}+ \hat{\rho}_{\tc{d},0}\hat{U}_I^{(2)^\dagger}\right) + \mathcal{O}(\lambda^4),
\end{align}
where
\begin{align}
    \hat{U}_I^{(1)} &= -\ii \int \dd V \, \hat{h}_\text{eff}(\mf x),\\
    \hat{U}_I^{(2)} &= -\int \dd V \dd V' \, \hat{h}_\text{eff}(\mf x)\hat{h}_\text{eff}(\mf x') \theta(t-t'),
\end{align}
with
\begin{equation}
    \hat{h}_\text{eff}(\mf x) = \lambda \hat{Q}_1^{\tc{a}}(\mf x) \hat{\phi}(\mf x) +\lambda \hat{Q}_1^{\tc{b}}(\mf x) \hat{\phi}(\mf x)  = \lambda(\Lambda_\tc{a}(\mf x) \hat{a}^{\tc{a}}_1+\Lambda^*_\tc{a}(\mf x) \hat{a}^{\tc{a}\dagger}_1 + \Lambda_\tc{b}(\mf x) \hat{a}^{\tc{b}}_1+\Lambda^*_\tc{b}(\mf x) \hat{a}^{\tc{b}\dagger}_1)\hat{\phi}(\mf x).
\end{equation}
This is exactly the leading order result when one considers the interaction of two harmonic oscillators interacting with a quantum field $\hat{\phi}(\mf x)$. That is, the final state of the modes can be written as
\begin{equation}
    \hat{\rho}_\tc{d} = \tr_\phi(\hat{U}_I (\hat{\rho}_{\tc{d},0} \otimes \hat{\rho}_\phi )\hat{U}_I^\dagger) + \mathcal{O}(\lambda^4), \quad \hat{U}_I = \mathcal{T}\exp\left(-\ii \int \dd V \, \hat{h}_\text{eff}(\mf x)\right).
\end{equation}
The leading order computations can then be carried on analogously to harmonic oscillator particle detectors (for details on this calculation, see e.g.~\cite{EricksonZero}). 



\chapter{The retarded propagator of the gravitational field}\label{app:retGrav}


The retarded propagator $G_R^{\mu\nu}{}_{\alpha'\beta'}(\mf x,\mf x')$ can be written as
\begin{align}
    G_R^{\mu\nu}{}_{\alpha'\beta'}(\mf x,\mf x') &= -\frac{1}{2\pi}\theta(t-t') \delta\left((t-t')^2 - |\bm x - \bm x'|^2\right)\mathcal{P}^{\mu\nu}{}_{\alpha'\beta'} \\
    &= -\frac{1}{4\pi | \bm x - \bm x'|} \delta(t-t' - |\bm x  - \bm x'|)\mathcal{P}^{\mu\nu}{}_{\alpha'\beta'},
\end{align}
where $\mathcal{P}$ is a bitensor. We then have the linearized metric given by $g_{\mu\nu} = \eta_{\mu\nu} + \sqrt{8\pi G}h_{\mu\nu}$, where
\begin{equation}
    h^{\mu\nu}(\mf x) = -\sqrt{8 \pi G}\int \dd V' G_R^{\mu\nu}{}_{\alpha' \beta'}(\mf x,\mf x') T^{\alpha' \beta'}(\mf x'),
\end{equation}
where $T_{\alpha\beta}$ denotes the stress-energy tensor of the source. For the case of a pointlike particle undergoing a trajectory $\mf z_1(t)$  with four-velocity $u_1^\mu(t)$, it reads
\begin{equation}
    T_1^{\mu\nu}(\mf x) = m_1 u_1^\mu(t) u_1^\nu(t) \frac{\delta^{(3)}(\bm x - \bm z_1(t))}{u_1^0(t) \sqrt{-g}},
\end{equation}
so that we obtain{\footnotesize
\begin{align}
    \int \dd V' G_R^{\mu\nu}{}_{\alpha'\beta'}(\mf x,\mf x')T_1^{\alpha' \beta'}(\mf x') & = -\frac{m_1}{4\pi} \int \dd t'  \frac{1}{u_1^0(t') |\bm x - \bm z_1(t')|} \delta(t-t' - |\bm x  - \bm z_1(t')|) \mathcal{P}^{\mu\nu}{}_{\alpha'\beta'}u_1^{\alpha'}(t') u_1^{\beta'}(t').
\end{align}}
Now, let $t_r$ be the retarded time such that $t - t_r - |\bm x - \bm z(t_r)| = 0$, so that
\begin{equation}
    \delta(t-t' - |\bm x  - \bm z_1(t')|) = \frac{\delta(t'-t_r)}{1- \hat{\bm r}_1(t_r)\!\cdot\! \dot{\bm z}_1(t_r)},
\end{equation}
where $\hat{\bm r_1}(t) = (\bm x  - \bm z_1(t))/|\bm x  - \bm z_1(t)|$ and we obtain
\begin{align}
    h_{(1)}^{\mu\nu}(\mf x) =- \sqrt{8 \pi G}\int \dd V' G_R^{\mu\nu}{}_{\alpha'\beta'}(\mf x,\mf x')T_1^{\alpha' \beta'}(\mf x') = \frac{ m_1}{4\pi}  \frac{\sqrt{8 \pi G}\mathcal{P}^{\mu\nu}{}_{\alpha'\beta'}u_1^{\alpha'}(t_r) u_1^{\beta'}(t_r)}{u_1^0(t_r)(1- \hat{\bm r}_1(t_r)\!\cdot\! \dot{\bm z}_1(t_r)) |\bm x - \bm z_1(t_r)|}.
\end{align}
The interaction Hamiltonian of a particle labelled $2$ with the gravitational potential sourced by particle $1$ will then be
\begin{align}
    {H}_{12}(t)  = - \frac{1}{2} \sqrt{8\pi G} \int \dd^3 {\bm x}\, h^{(1)}_{\mu\nu}(\mf x) T_2^{\mu\nu}(\mf x) &=- \frac{G m_1 m_2}{|\bm z_2(t) - \bm z_1(t_r)|} \frac{\mathcal{P}_{\mu\nu}{}_{\alpha'\beta'}u_2^{\mu}(t) u_2^{\nu}(t)u_1^{\alpha'}(t_r) u_1^{\beta'}(t_r)}{u_2^0(t)u_1^0(t_{r_{12}})(1- \hat{\bm r}_{12}\!\cdot\! \dot{\bm z}_1(t_{r_{12}}))},\nonumber\\
    &=- \frac{G m_1 m_2}{|\bm z_2(t) - \bm z_1(t_{r_{12}})|} \frac{2(\eta_{\mu\nu}u_2^{\mu}(t) u_1^{\nu}(t_{r_{12}}))^2-1}{u_2^0(t)u_1^0(t_{r_{12}})(1- \hat{\bm r}_{12}\!\cdot\! \dot{\bm z}_1(t_{r_{12}}))},
\end{align}
where $t_{r_{12}}$ is the solution to $t - t_{r_{12}} = |\bm z_2(t) - \bm z_1(t_{r_{12}})|$ and $\hat{\bm r}_{12} = (\bm z_2(t)  - \bm z_1(t_{r_{12}}))/|\bm z_2(t)  - \bm z_1(t_{r_{12}})|$, and we used $\mathcal{P}_{\mu\nu}{}_{\alpha'\beta'} = \eta_{\mu\alpha'}\eta_{\nu \beta'} + \eta_{\mu \beta'} \eta_{\nu \alpha'} - \eta_{\mu\nu} \eta_{\alpha'\beta'}$. The total interaction between the two particles is then given by
\begin{align}
    H_I(t) = \frac{1}{2}\left(H_{12}(t) + H_{21}(t)\right) =&- \frac{G m_1 m_2}{|\bm z_2(t) - \bm z_1(t_{r_{12}})|}\frac{(\eta_{\mu\nu}u_2^{\mu}(t) u_1^{\nu}(t_{r_{12}}))^2-1/2}{u_2^0(t)u_1^0(t_{r_{12}})(1- \hat{\bm r}_{12}\!\cdot\! \dot{\bm z}_1(t_{r_{12}}))}\\&- \frac{G m_1 m_2}{|\bm z_1(t) - \bm z_2(t_{r_{21}})|} \frac{(\eta_{\mu\nu}u_1^{\mu}(t) u_2^{\nu}(t_{r_{21}}))^2-1/2}{u_1^0(t)u_2^0(t_{r_{21}})(1- \hat{\bm r}_{12}\!\cdot\! \dot{\bm z}_2(t_{r_{21}}))},
\end{align}
$t_{r_{21}}$ is the solution to $t - t_{r_{21}} = |\bm z_2(t) - \bm z_1(t_{r_{21}})|$ and $\hat{\bm r}_{21} = (\bm z_1(t)  - \bm z_2(t_{r_{21}}))/|\bm z_1(t)  - \bm z_2(t_{r_{21}})|$, and we used $\mathcal{P}_{\mu\nu}{}_{\alpha'\beta'} = \eta_{\mu\alpha'}\eta_{\nu \beta'} + \eta_{\mu \beta'} \eta_{\nu \alpha'} - \eta_{\mu\nu} \eta_{\alpha'\beta'}$.
Notice that in the non-relativistic limit, we have $t_{r_{12}}\approx t_{r_{21}} \approx t$, $u_1^0(t) \approx u_2^0(t) \approx 1$ and $\eta_{\mu\nu} u_1^\mu(t)u_2^\nu \approx 1$, allowing one to recover the non-local Newtonian interaction between the particles:
\begin{equation}
    H_I(t)\approx - \frac{G m_1 m_2}{|\bm z_1(t) - \bm z_2(t)|}.
\end{equation}


Also notice that
\begin{align}
    \int \dd t\, H_I(t) &= 2 \pi G\int \dd V \dd V'\left( T^1_{\mu\nu}(\mf x) G_R^{\mu\nu}{}_{\alpha'\beta'}(\mf x, \mf x')T_2^{\alpha' \beta'}(\mf x') +  T^2_{\mu\nu}(\mf x) G_R^{\mu\nu}{}_{\alpha'\beta'}(\mf x, \mf x')T_1^{\alpha' \beta'}(\mf x')\right) \nonumber\\
    &= 2 \pi G\int \dd V \dd V'\left( T^1_{\mu\nu}(\mf x) G_R^{\mu\nu}{}_{\alpha'\beta'}(\mf x, \mf x')T_2^{\alpha' \beta'}(\mf x') +  T^2_{\alpha'\beta'}(\mf x') G_R^{\alpha'\beta'}{}_{\mu\nu}(\mf x', \mf x)T_1^{\mu \nu}(\mf x')\right) \nonumber\\
    &= 2 \pi G\int \dd V \dd V' T^1_{\mu\nu}(\mf x) \Delta^{\mu\nu}{}_{\alpha'\beta'}(\mf x, \mf x')T_2^{\alpha' \beta'}(\mf x'),
\end{align}
where $\Delta^{\mu\nu\alpha'\beta'}(\mf x, \mf x') = G_R^{\mu\nu\alpha'\beta'}(\mf x, \mf x') + G_A^{\mu\nu\alpha'\beta'}(\mf x,\mf x')$ and we used $G_R^{\alpha'\beta'\mu\nu}(\mf x',\mf x) = G_A^{\mu\nu\alpha'\beta'}(\mf x,\mf x')$.




\chapter{Final states in the Quantum Field Approach to Gravity Mediated Entanglement}\label{app:FinalStates}

The initial state of the two particles in matrix form in the basis $\{\ket{L_1L_2},\ket{R_1 L_2},\ket{L_1R_2},\ket{R_1R_2}\}$ reads
\begin{equation}
    \hat{\rho} = \frac{1}{4}
    \begin{pmatrix}
        1 & 1 & 1 & 1\\
        1 & 1 & 1 & 1\\
        1 & 1 & 1 & 1\\
        1 & 1 & 1 & 1        
    \end{pmatrix}.
\end{equation}
To leading order in the gravitational field, the updated state of the particles evolving with respect to the qc-interaction can be written as $\hat{\rho} +  \delta\hat{\rho}_c$, where
\begin{equation}
    \delta\hat{\rho}_c = -\frac{\ii \lambda^2}{4}
    \begin{pmatrix}
        0 & \Delta_{R_1L_2}-\Delta_{L_1L_2} &\Delta_{L_1R_2}-\Delta_{L_1L_2} & \Delta_{R_1R_2}-\Delta_{L_1L_2}\\
        -\Delta_{R_1L_2}+\Delta_{L_1L_2} & 0 & \Delta_{L_1R_2}-\Delta_{R_1L_2} & \Delta_{R_1R_2}-\Delta_{R_1L_2}\\
       -\Delta_{L_1L_2}+\Delta_{L_1R_2} & -\Delta_{L_1R_2}+\Delta_{R_1L_2} & 0 & \Delta_{L_1R_2}-\Delta_{R_1R_2}\\
        -\Delta_{L_1L_2}+\Delta_{R_1R_2} & -\Delta_{R_1L_2}+\Delta_{R_1R_2} & -\Delta_{L_1R_2}+\Delta_{R_1R_2} & 0      
    \end{pmatrix}.
\end{equation}
In the quantum case, the final state of the particles to leading order can be written as $\hat{\rho} + (\delta\hat{\rho}_c+\delta\hat{\rho}_q+\delta \hat{\rho}_l)$, where
\begin{equation}
    \delta\hat{\rho}_l = \frac{\lambda^2}{2} (\mathcal{L}_\textsc{v}-\mathcal{L}_{\textsc{i}})
    \begin{pmatrix}
        0  &  1 & 1 & 2\\
        1  &  0 & 2 & 1\\
        1  &  2 & 0 & 1\\
        2  &  1 & 1 & 0
    \end{pmatrix},
\end{equation}
and
\begin{equation}
    \delta\hat{\rho}_q = \frac{\lambda^2}{4}(H_{L_1R_2} \!\!+ \!\!H_{R_1L_2} \!\!-\!\! H_{L_1L_2}\!\!-\!\!H_{R_1R_2})
    \begin{pmatrix}
        0 & 0 & 0 & 1\\
        0 & 0 & -1 & 0\\
        0 & -1 & 0 & 0\\
        1 & 0 & 0 & 0
    \end{pmatrix},
\end{equation}
where $H_{p_1p_2}$ denotes the integrated Hadamard function along each pair of paths for particles 1 and 2,
\begin{equation}
    H_{p_1p_2} = \int \dd V \dd V' T^{\mu\nu}_{p_1}H_{\mu\nu\alpha'\beta'}(\mf x,\mf x')T_{p_2}^{\alpha'\beta'}(\mf x').
\end{equation}
while $\mathcal{L}_\textsc{v}$ and $\mathcal{L}_{\textsc{i}}$ are noise terms due to the interaction of each particle with the vacuum of the field. These are explicitly given by
\begin{align}
    \mathcal{L}_{\textsc{v}} &= \int \dd V \dd V' T^{\mu\nu}_{L_i}\langle \hat{h}_{\mu\nu}(\mf x) \hat{h}_{\alpha'\beta'}(\mf x')\rangle_0 T_{L_i}^{\alpha'\beta'}(\mf x')= \int \dd V \dd V' T^{\mu\nu}_{R_i}\langle \hat{h}_{\mu\nu}(\mf x) \hat{h}_{\alpha'\beta'}(\mf x')\rangle_0 T_{R_i}^{\alpha'\beta'}(\mf x'),\\
    \mathcal{L}_{\textsc{i}} &= \int \dd V \dd V' T^{\mu\nu}_{L_i}\langle \hat{h}_{\mu\nu}(\mf x) \hat{h}_{\alpha'\beta'}(\mf x')\rangle_0 T_{R_i}^{\alpha'\beta'}(\mf x'),
\end{align}
where, due to the choice of paths, the indices $i$ can be $1$ or $2$, and still yield the same result due to the fact that the paths are related to translations and rotations, which are symmetries of the quantum field theory. Then, the $\mathcal{L}_{\textsc{v}}$ term is a local noise associated to each path, while the $\mathcal{L}_{\textsc{i}}$ term represents an interference term associated with each particle undergoing the superposition of paths. Also notice that due to the fact that the propagator decreases with distance, we have $\mathcal{L}_{\textsc{i}}\leq \mathcal{L}_{\textsc{v}}$, with equality holding only if the paths $1$ and $2$ are identical. Moreover, these vacuum noise terms decay fast with the interaction time, so they are negligible for long interaction times. Then we can interpret $\delta{\hat{\rho}}_l$ as a local vacuum contribution, and $\delta\hat{\rho}_q$ can be seen as the additional correlation contribution due to the quantum nature of the field. 

